\section{Problem Statement}
\label{sec:problem}

\subsection{Current Limitations}
Existing retail forecasting tooling accessible to non-specialist organizations typically exhibits one or more of the following limitations. First, many tools commit to a single model family, such as a single statistical method or a single machine learning regressor, and therefore perform inconsistently across the heterogeneous demand patterns present even within one retailer's catalogue, a pattern well documented by cross-method forecasting competitions~\cite{makridakis2020m4,makridakis2022m5}. Second, tools that do support multiple models frequently select among them using only a single point-accuracy metric, most commonly root-mean-squared error, which is sensitive to the scale of the series and does not account for computation cost, forecast stability, or the risk that one poorly calibrated model can distort a comparison against otherwise well-fitted competitors. Third, deep learning models are often applied uniformly regardless of dataset size, despite established evidence that they require substantially more observations than statistical or tree-based alternatives to generalize reliably~\cite{lim2021survey}; on short retail series this risks reporting an overfit model as though it were reliable. Fourth, and most consequential for adoption by non-technical stakeholders, the forecast is typically delivered as a bare number or chart with no accompanying explanation of model choice, uncertainty, or business implication, which limits the degree to which a retail manager can trust and act on the output.

\subsection{Research Gap}
The literature reviewed in Section~\ref{sec:literature} addresses these limitations largely in isolation: forecasting-competition research establishes which model families perform well on which data characteristics, evaluation-methodology research establishes how models should be validated, and explainable-AI research establishes how a trained model's behavior can be attributed or narrated. There is comparatively little published work describing a single, integrated system that combines a broad multi-family model ensemble, a robust and auditable multi-criteria selection rule, an explicit data-driven gate on which models are even attempted, and a language-model explanation layer that is architecturally constrained to remain factually consistent with the system's own deterministic computations. This integration gap, rather than any single algorithmic deficiency, is the specific research gap this paper addresses.

\subsection{Objectives}
The system described in this paper is designed to satisfy the following objectives.
\begin{enumerate}
    \item Accept an arbitrary time-indexed dataset in common tabular file formats and automatically identify the date column, the numeric target column, and any usable exogenous or grouping columns, without requiring the uploader to pre-process the file.
    \item Train a broad ensemble of forecasting models spanning statistical, classical machine learning, deep learning, and intermittent-demand families, evaluated identically via walk-forward cross-validation.
    \item Select the best-performing model using a composite criterion that is robust to a single outlier model and that reflects accuracy, stability, and computational cost jointly, rather than a single error metric alone.
    \item Withhold computationally expensive deep learning models from datasets below a data-driven minimum size, and communicate this decision transparently rather than silently degrading forecast quality.
    \item Incorporate real exogenous business features, specifically price and promotional indicators, into the subset of models capable of using them, when such columns are present in the uploaded dataset.
    \item Provide a natural-language explanation layer, grounded in retrieved domain knowledge and in the pipeline's own deterministic outputs, that answers user questions and produces a proactive business report without contradicting the system's own computed metrics or labels.
\end{enumerate}

\subsection{Expected Improvements}
Relative to a single-model or single-metric baseline, the proposed system is expected to improve robustness across heterogeneous retail series by virtue of its broad ensemble and adaptive model gating, to improve the trustworthiness of reported uncertainty by deriving forecast intervals from each model's own historical error rather than an arbitrary fixed percentage, and to improve the practical usability of the forecast for non-technical stakeholders through grounded natural-language explanation. These improvements are evaluated empirically in Sections~\ref{sec:experimental} and~\ref{sec:results}.
